\clearpage\thispagestyle{empty}
\pdfbookmark[0]{Quellen zu Meridiankreis (5. Juni 1820)}{meridianjune1820-sources}
\begin{center}\textbf{Quellen zu Meridiankreis (5. Juni 1820)}\end{center}
\noindent Diese KI-Übertragung ins heutige Deutsch folgt \textit{Meridiankreis} in den \href{https://gdz.sub.uni-goettingen.de/id/PPN236007467}{\textit{Werken}, Band~VI} (Göttingen, 1874), S.~429--433. Der \href{https://images.p.sub.uni-goettingen.de/iiif/image/gdz:PPN319721507_1820:00000901/full/1600,/0/default.jpg}{Erstdruck} erschien unter \textit{Göttingen} in den \textit{Göttingischen gelehrten Anzeigen}, Nr.~91, 5. Juni 1820, S.~905--912. Das digitalisierte Exemplar stammt aus der Niedersächsischen Staats- und Universitätsbibliothek Göttingen. Der Bericht verwendet die erste Person Plural und nennt Gauß in der dritten Person als Beobachter. Der Verweis auf Reichenbachs Passageinstrument führt zum \hyperlink{native-427-part-lower}{Bericht vom Oktober 1819}.
\par\smallskip\noindent Der Erstdruck gibt auf \href{https://images.p.sub.uni-goettingen.de/iiif/image/gdz:PPN319721507_1820:00000902/full/1600,/0/default.jpg}{S.~906} 14 Sekunden an. Die Werkausgabe druckt vor der 14 einen Punkt. Auf \href{https://images.p.sub.uni-goettingen.de/iiif/image/gdz:PPN319721507_1820:00000907/full/1600,/0/default.jpg}{S.~911} nennt der Erstdruck in der Überschrift zur Zenitdistanz des Polarsterns den Kollimationsfehler. Die Werkausgabe nennt dort den Kulminationsfehler. In der Zeitangabe zum Mars am 11. April steht in der Werkausgabe zwischen 20 und 3 ein Zwischenraum, im \href{https://images.p.sub.uni-goettingen.de/iiif/image/gdz:PPN319721507_1820:00000908/full/1600,/0/default.jpg}{Erstdruck} ein Komma. Die Übertragung folgt der Werkausgabe. KI-Übertragung und KI-Satz.
\clearpage
